\documentclass{ximera}
%verbeterd 30/10/26
\title{Samenvatting gemengde schakelingen}
%\author{Daan Lipkens}
\begin{document}
\begin{abstract}
\end{abstract}
\begin{onlineOnly}
\maketitle
\end{onlineOnly}
% pagestyle aanpassen zodat er meer op 1 blad kan
\newgeometry{left=3cm,bottom=0.1cm}
\pagestyle{empty}

\begin{conclusion}[title={Samenvatting}]\nl

    \samenbox{Grootheden en wet van Ohm}{
        De stroomsterkte $I = \frac{\Delta Q}{\Delta t}$ meet je in ampère ($\mathrm{A}$), de spanning $U$ in volt ($\mathrm{V}$) en de weerstand $R$ in ohm ($\mathrm{\Omega}$). Voor een ohmse weerstand geldt de wet van Ohm:
        \begin{equation*}
            U = R \cdot I
        \end{equation*}
        Een ampèremeter schakel je in serie, een voltmeter parallel.
    }

    \samenbox{Serieschakeling}{
        Eén pad, geen vertakkingen.
        \begin{itemize}
            \item Stroom overal gelijk: $I = I_1 = I_2 = \dots$
            \item Spanningen tellen op: $U = U_1 + U_2 + \dots$
            \item Vervangingsweerstand: $R_{\text{s}} = R_1 + R_2 + \dots$
        \end{itemize}
        De grootste weerstand krijgt de grootste spanning.
    }

    \samenbox{Parallelschakeling}{
        Meerdere paden met vertakkingspunten.
        \begin{itemize}
            \item Spanning over elke tak gelijk: $U = U_1 = U_2 = \dots$
            \item Stromen tellen op: $I = I_1 + I_2 + \dots$
            \item Vervangingsweerstand: $\frac{1}{R_{\text{p}}} = \frac{1}{R_1} + \frac{1}{R_2} + \dots$
        \end{itemize}
        $R_{\text{p}}$ is kleiner dan de kleinste weerstand. De kleinste weerstand krijgt de grootste stroom.
    }

    \samenbox{Gemengde schakeling}{
        \begin{enumerate}
            \item Vervang groepen in serie of parallel door één vervangingsweerstand, tot er $R_{\text{tot}}$ overblijft.
            \item Bereken $I = \frac{U}{R_{\text{tot}}}$.
            \item Werk terug: in serie dezelfde stroom, in parallel dezelfde spanning.
            \item Controleer met de som van spanningen of stromen.
        \end{enumerate}
    }

\end{conclusion}

% restore pagestyle
\pagestyle{plain}
\restoregeometry
\end{document}